\begin{lemma}[Smooth approximation of a compactly supported $W^{1,q}$ function]
Let $n\ge1$ be an integer and $q\ge1$ real, let $\Omega\subseteq\R^n$ be open,
and let $\varphi\in W^{1,q}(\Omega)$ with weak gradient $g_\varphi$
(\noderef{MemW1pOn}) have compact support with
$\operatorname{supp}\varphi\subseteq\Omega$.  Then there are a compact set
$K\subseteq\Omega$ with $\operatorname{supp}\varphi\subseteq K$ and functions
$\Phi_m\in C^\infty(\R^n)$, $m\in\mathbb{N}$, each with compact support contained in
$K$, such that
\[
\int_K|\Phi_m-\varphi|^q\,dx\longrightarrow0
\qquad\text{and}\qquad
\int_K\bigl|D_i\Phi_m-(g_\varphi)_i\bigr|^q\,dx\longrightarrow0
\quad(m\to\infty)
\]
for each $i\in\{1,\dots,n\}$.  All powers are real powers of nonnegative reals.
\end{lemma}
\begin{proof}
Write $K_0=\operatorname{supp}\varphi$.

\emph{The degenerate case.}  If $K_0=\emptyset$ then $\varphi$ vanishes
identically, and we may take $K=\emptyset$ --- compact, contained in $\Omega$,
and containing $K_0$ --- together with $\Phi_m\equiv0$, which is $C^\infty$ with
compact support $\emptyset\subseteq K$.  Every integral over $\emptyset$ is $0$,
so both asserted limits are the constant sequence $0$.  Assume from now on
$K_0\ne\emptyset$.

\emph{Step 1: geometry.}  The set $K_0$ is compact, $\R^n\setminus\Omega$ is
closed, and the two are disjoint, so
$d=\operatorname{dist}(K_0,\R^n\setminus\Omega)>0$ (with $d=+\infty$ when
$\Omega=\R^n$).  Fix $\delta_0>0$ with $3\delta_0<d$ and for $s>0$ put
\[
K_s'=\{x\in\R^n:\operatorname{dist}(x,K_0)\le s\} .
\]
Each $K_s'$ is closed, the distance function to $K_0$ being continuous, and
bounded because $K_0$ is bounded, hence compact; and $K_s'\subseteq\Omega$
whenever $s\le3\delta_0$, since a point outside $\Omega$ has distance at least
$d>3\delta_0$ from $K_0$.  Set
\[
K=K_{2\delta_0}' ,
\]
so $K$ is compact, $K\subseteq\Omega$ and $K_0\subseteq K$, which are the first
three assertions about $K$.

\emph{Step 2: the data as functions on $\R^n$.}  By \noderef{MemW1pOn} we have
$\varphi\in L^q(\Omega)$ and $|g_\varphi|\in L^q(\Omega)$, both almost
everywhere strongly measurable on $\Omega$.
The function $\varphi$ vanishes at every point outside
$K_0=\operatorname{supp}\varphi$, by the definition of the support; in
particular $\varphi$ agrees with $\varphi\chi_{K_0}$ everywhere, so, being
almost everywhere strongly measurable on $\Omega$ and identically $0$ on the
complement of $K_0\subseteq\Omega$, it is almost everywhere strongly measurable
on all of $\R^n$.  Hence $\varphi\in L^q(\R^n)$ with
$\int_{\R^n}|\varphi|^q=\int_{K_0}|\varphi|^q\le\int_\Omega|\varphi|^q<\infty$.
Put
\[
\check g=g_\varphi\,\chi_\Omega ,
\]
the extension of $g_\varphi$ by $0$ off the (measurable, indeed open) set
$\Omega$; then $\check g$ is almost everywhere strongly measurable on $\R^n$ ---
almost everywhere strongly measurable on $\Omega$ by \noderef{MemW1pOn} and
identically $0$ off $\Omega$ --- with $\check g=g_\varphi$ on
$\Omega$, and $|\check g|\in L^q(\R^n)$ with
$\int_{\R^n}|\check g|^q=\int_\Omega|g_\varphi|^q<\infty$.  Each component
satisfies $|\check g_i|\le|\check g|$, so $\check g_i\in L^q(\R^n)$ as well.
Finally, each of $\varphi$ and $\check g_i$ is locally integrable on $\R^n$: on
a compact $K'$ we have $|K'|<\infty$ and $q\ge1$, so H\"older on a set of finite
measure (Preamble) gives $L^q(K')\subseteq L^1(K')$.

\emph{Step 3: the mollifier.}  For each $\delta>0$ let $\rho_\delta$ be the mollifier at scale $\delta$
supplied by the Preamble item on a standard mollifier:
\[
\rho_\delta\in C^\infty(\R^n),\qquad \rho_\delta\ge0,\qquad
\operatorname{supp}\rho_\delta\subseteq\overline{B_\delta(0)}\ \text{compact},
\qquad \int_{\R^n}\rho_\delta\,dz=1 .
\]
Those four properties are the only ones used below.  In particular no explicit
formula for $\rho_\delta$ is needed --- the Preamble item supplies the family
directly at every scale, so $\rho_\delta$ is not obtained here by rescaling a
single kernel and no change of variables has to be justified.  For $w$ locally integrable on $\R^n$ write
$(w*\rho_\delta)(x)=\int_{\R^n}w(y)\rho_\delta(x-y)\,dy$.  Choose
$\delta_m=\delta_0/(m+1)$ for $m\in\mathbb{N}$, so $0<\delta_m\le\delta_0$ and
$\delta_m\downarrow0$, and define
\[
\Phi_m=\varphi*\rho_{\delta_m} .
\]

\emph{Step 4: $\Phi_m$ is smooth and supported in $K$.}  By the Preamble item on
convolution with a smooth compactly supported kernel, applied to the locally
integrable $\varphi$ of Step 2 and the kernel $\rho_{\delta_m}$, each $\Phi_m$
is of class $C^\infty$ on $\R^n$ and
\[
D_i\Phi_m=\varphi*(D_i\rho_{\delta_m})\qquad(i=1,\dots,n).
\tag{1}
\]
For the support: if $\Phi_m(x)\ne0$ then the integrand
$y\mapsto\varphi(y)\rho_{\delta_m}(x-y)$ does not vanish identically, so there
is $y$ with $\varphi(y)\ne0$ and $|x-y|\le\delta_m$; then $y\in K_0$ and
$\operatorname{dist}(x,K_0)\le\delta_m\le\delta_0$.  Hence
$\{\Phi_m\ne0\}\subseteq K_{\delta_0}'$, and $K_{\delta_0}'$ being closed, also
$\operatorname{supp}\Phi_m\subseteq K_{\delta_0}'\subseteq
K_{2\delta_0}'=K$.  In particular $\Phi_m$ has compact support, being a closed
subset of the compact $K$.

\emph{Step 5: the derivative of $\Phi_m$ on $K$ is the mollified weak gradient.}
Let $x\in K$, let $m\in\mathbb{N}$ and let $i$ be an index.  If $|y-x|\le\delta_m$ then
\[
\operatorname{dist}(y,K_0)\le|y-x|+\operatorname{dist}(x,K_0)
\le\delta_m+2\delta_0\le3\delta_0 ,
\]
so $\overline{B_{\delta_m}(x)}\subseteq K_{3\delta_0}'\subseteq\Omega$.  Define
$\psi(y)=\rho_{\delta_m}(x-y)$.  Then $\psi\in C^\infty(\R^n)$ and
$\operatorname{supp}\psi\subseteq\overline{B_{\delta_m}(x)}$, which is compact
and contained in $\Omega$ by the display; so $\psi$ belongs to the test class of
\noderef{HasWeakGradientOn} for $\Omega$.  Differentiating in the $y$ variable,
$D_i\psi(y)=-(D_i\rho_{\delta_m})(x-y)$.  The defining identity of
\noderef{HasWeakGradientOn}, which the pair $(\varphi,g_\varphi)$ satisfies on
$\Omega$ by \noderef{MemW1pOn}, applied to $\psi$ and the index $i$, reads
$\int_\Omega\varphi\,D_i\psi\,dy=-\int_\Omega(g_\varphi)_i\,\psi\,dy$, that is
\[
-\int_\Omega\varphi(y)(D_i\rho_{\delta_m})(x-y)\,dy
=-\int_\Omega(g_\varphi)_i(y)\rho_{\delta_m}(x-y)\,dy .
\]
Both integrands vanish off $\overline{B_{\delta_m}(x)}\subseteq\Omega$, where
$\check g=g_\varphi$, so each integral over $\Omega$ equals the corresponding
integral over $\R^n$; comparing with (1) and multiplying by $-1$,
\[
D_i\Phi_m(x)=\bigl(\check g_i*\rho_{\delta_m}\bigr)(x)
\qquad(x\in K,\ i=1,\dots,n,\ m\in\mathbb{N}).
\tag{2}
\]

\emph{Step 6: the $L^q$ contraction of convolution against a probability
density.}  Let $w\in L^q(\R^n)$ and $\delta>0$.  For each fixed $x$ the measure
$\mu_x(dy)=\rho_\delta(x-y)\,dy$ is a nonnegative measure of total mass
$\mu_x(\R^n)=\int\rho_\delta(x-y)\,dy=1$, by translation and reflection
invariance of Lebesgue measure (Preamble) and $\int\rho_\delta=1$.  H\"older on
a set of finite measure (Preamble), applied on $\R^n$ with the measure $\mu_x$
and the exponent $q$, gives
\[
|(w*\rho_\delta)(x)|\le\int_{\R^n}|w|\,d\mu_x
\le\mu_x(\R^n)^{1-\frac1q}\Bigl(\int_{\R^n}|w|^q\,d\mu_x\Bigr)^{1/q}
=\Bigl(\int_{\R^n}|w(y)|^q\rho_\delta(x-y)\,dy\Bigr)^{1/q},
\]
so, raising to the power $q$ (which is nondecreasing on $[0,\infty)$) and
integrating in $x$,
\[
\int_{\R^n}|w*\rho_\delta|^q\,dx
\le\int_{\R^n}\!\!\int_{\R^n}|w(y)|^q\rho_\delta(x-y)\,dy\,dx
=\int_{\R^n}|w(y)|^q\Bigl(\int_{\R^n}\rho_\delta(x-y)\,dx\Bigr)dy
=\int_{\R^n}|w|^q\,dy ,
\]
the exchange by Tonelli (Preamble), the integrand
$(x,y)\mapsto|w(y)|^q\rho_\delta(x-y)$ being nonnegative and measurable as the
product of a measurable function of $y$ with a continuous function of $(x,y)$,
and the inner integral being $1$ as above.  Thus
\[
\|w*\rho_\delta\|_{L^q(\R^n)}\le\|w\|_{L^q(\R^n)} .
\tag{3}
\]

\emph{Step 7: mollifications converge in $L^q$.}  We claim that for every
$w\in L^q(\R^n)$,
\[
\|w*\rho_\delta-w\|_{L^q(\R^n)}\longrightarrow0\qquad(\delta\downarrow0).
\tag{4}
\]
Let $\sigma>0$.  By the Preamble item on density of the compactly supported
continuous functions in $L^q$, applied at the exponent $q$, there is a
continuous $v$ with compact support and $\|w-v\|_{L^q}\le\sigma/3$.  By the
$L^q$ triangle inequality (Minkowski, Preamble) and (3),
\[
\|w*\rho_\delta-w\|_{L^q}
\le\|(w-v)*\rho_\delta\|_{L^q}+\|v*\rho_\delta-v\|_{L^q}+\|v-w\|_{L^q}
\le2\|w-v\|_{L^q}+\|v*\rho_\delta-v\|_{L^q} .
\]
For the middle term, $v$ is continuous with compact support, hence bounded and
uniformly continuous, and since $\rho_\delta$ is a probability density supported
in $\overline{B_\delta(0)}$,
\[
|(v*\rho_\delta)(x)-v(x)|
=\Bigl|\int_{\R^n}\bigl(v(x-z)-v(x)\bigr)\rho_\delta(z)\,dz\Bigr|
\le\sup_{|z|\le\delta}\ \sup_{x\in\R^n}|v(x-z)-v(x)|
\xrightarrow[\delta\downarrow0]{}0
\]
by uniform continuity; and for $\delta\le1$ every function
$v*\rho_\delta-v$ vanishes off the fixed compact set
$C=\{x:\operatorname{dist}(x,\operatorname{supp}v)\le1\}$, so
$\|v*\rho_\delta-v\|_{L^q}^q\le|C|\cdot\sup_x|(v*\rho_\delta)(x)-v(x)|^q\to0$.
Hence $\limsup_{\delta\downarrow0}\|w*\rho_\delta-w\|_{L^q}\le2\sigma/3<\sigma$,
and as $\sigma>0$ was arbitrary, (4) holds.

\emph{Step 8: conclusion.}  Apply (4) to $w=\varphi$, which lies in
$L^q(\R^n)$ by Step 2, along the sequence $\delta_m\downarrow0$: since
$\Phi_m=\varphi*\rho_{\delta_m}$ and the integral over $K$ of a nonnegative
function is at most its integral over $\R^n$,
\[
\int_K|\Phi_m-\varphi|^q\,dx\le\int_{\R^n}|\varphi*\rho_{\delta_m}-\varphi|^q\,dx
=\|\varphi*\rho_{\delta_m}-\varphi\|_{L^q(\R^n)}^q\longrightarrow0 ,
\]
which is the first asserted limit.  Fix an index $i$ and apply (4) to
$w=\check g_i\in L^q(\R^n)$: by (2) we have
$D_i\Phi_m=\check g_i*\rho_{\delta_m}$ on $K$, and $\check g_i=(g_\varphi)_i$ on
$K\subseteq\Omega$, so
\[
\int_K\bigl|D_i\Phi_m-(g_\varphi)_i\bigr|^q\,dx
=\int_K\bigl|\check g_i*\rho_{\delta_m}-\check g_i\bigr|^q\,dx
\le\bigl\|\check g_i*\rho_{\delta_m}-\check g_i\bigr\|_{L^q(\R^n)}^q
\longrightarrow0 ,
\]
which is the second asserted limit.  Together with Steps 1 and 4 --- $K$
compact, $K\subseteq\Omega$, $K_0\subseteq K$, each $\Phi_m$ of class
$C^\infty$ with compact support contained in $K$ --- this is the assertion.
\end{proof}
